\section{Introduction}
\label{sec:intro}

\subsection{The first row}
\label{ssec:firstrow}

Let
\begin{equation}\label{eq:apery-a}
  a_n \;=\; \sum_{k=0}^{n}\bin nk^2\bin{n+k}k^2
  \;=\;1,\;5,\;73,\;1445,\;33001,\;819005,\;21460825,\;\dots
\end{equation}
be Ap\'ery's numbers for $\zeta(3)$ (OEIS \texttt{A005259}). They satisfy a congruence of
Lucas type: for every prime $p$, every $a\ge 0$ and every $r$ with $0\le r<p$,
\begin{equation}\label{eq:arow}
  a_{ap+r}\;\equiv\;a_a\,a_r \pmod p ,
\end{equation}
and consequently $a_n\equiv\prod_i a_{n_i}\pmod p$ for the base-$p$ expansion
$n=\sum_i n_ip^i$. This is classical \cite{Gessel1982}, and it is by now known
to be a structural, not accidental, feature of the whole \emph{sporadic} class:
Malik and Straub \cite[Thm.~3.1]{MalikStraub} prove \eqref{eq:arow} for every one of the
fifteen sporadic Ap\'ery-like sequences in Zagier's classification, by McIntosh's
constant-term method, and Straub \cite{Straub2301} sharpens it to a congruence modulo
$p^2$ for all fifteen. Beyond Lucas congruences there is a supercongruence:
$a_{p^rm}\equiv a_{p^{r-1}m}\pmod{p^{3r}}$ for $p\ge 5$
\cite[Thm.~3.2]{Straub1401}, a statement whose history runs through Beukers' conjecture
and Coster's proof, as recorded in the survey of Osburn and Sahu \cite{OsburnSahu}. On the
crystalline side, Beukers and Vlasenko's Dwork-crystal papers
\cite{BVone,BVtwo,BVthree} and Mellit--Vlasenko \cite{MellitVlasenko} put the same
first-row congruences into a Frobenius framework, where they appear as statements about
the \emph{unit root} $F(t)/F(t^\sigma)$ of the holomorphic period.

In short: the arithmetic of the integral solution of an Ap\'ery-like recurrence is
well-trodden ground.

\subsection{The second row, and the gap}
\label{ssec:gap}

Ap\'ery's proof of the irrationality of $\zeta(3)$ uses a \emph{second} solution of the
same recurrence,
\begin{equation}\label{eq:apery-b}
  b_n \;=\; \sum_{k=0}^n \bin nk^2\bin{n+k}k^2
     \Bigl(\HH{3}{n} \;+\; \sum_{m=1}^{k}
       \frac{(-1)^{m-1}}{2m^3\bin nm\bin{n+m}m}\Bigr),
  \qquad \HH{r}{x}=\sum_{m=1}^{x}\frac1{m^r},
\end{equation}
with $b_n/a_n\to\zeta(3)$; the first values are
$0,\,6,\,\tfrac{351}{4},\,\tfrac{62531}{36},\,\tfrac{11424695}{288},\dots$
(OEIS \texttt{A059415} for the numerators). The pair $(a_n,b_n)$ is not integral: $b_n$
has denominators, bounded by $d_n^3b_n\in\Z$ with $d_n=\lcm(1,\dots,n)$.

For this second row, nothing of the shape \eqref{eq:arow} appears in the literature. We
surveyed the relevant technology in detail, working only from fetched sources, and the
picture is uniform:

\begin{itemize}[leftmargin=1.4em]
\item \emph{Lucas congruences for constant-term sequences} (Malik--Straub
  \cite{MalikStraub}, Henningsen--Straub \cite{HenningsenStraub}, Mellit--Vlasenko
  \cite{MellitVlasenko}, Straub \cite{Straub1401,Straub2301}) treat the integral solution
  only.
\item \emph{Dwork crystals} \cite{BVone,BVtwo,BVthree} construct the logarithmic second
  solution $F(t)\log t+G(t)$, whose coefficients are exactly our $B(n)$; they explain the
  $p^w$ normalisation, since the Cartier matrix on the rank-two quotient is
  $\Lambda=Y\Lambda_0(Y^\sigma)^{-1}$ with
  $\Lambda_0=\begin{psmallmatrix}\alpha_0&\alpha_1\\0&p\alpha_0\end{psmallmatrix}$, so
  that the second-solution Frobenius eigenvalue carries exactly one extra factor of $p$
  per Hodge step \cite[Prop.~7.7]{BVthree}; and they prove $p$-integrality of the mirror
  map $q(t)=t\exp(G/F)$ \cite[Cor.~7.11]{BVthree}. What they do \emph{not} prove is any
  congruence modulo $p$ or $p^s$ for the second solution: the corresponding deepening is
  their Conjecture~7.5, explicitly open, and stated for the \emph{unit root} at that.
\item \emph{Mirror-map integrality} (Delaygue--Rivoal--Roques \cite{DRR}, Delaygue
  \cite{Delaygue}, \cite{MirrorNote}) and \emph{strong Frobenius structures}
  (Vargas-Montoya \cite{VMcanon,VMmum}) control the denominators of the second-solution
  ratio, and in \cite{VMmum} prove Lucas congruences --- but explicitly only for the
  \emph{holomorphic} solution.
\item \emph{$p$-adic zeta values} sit in the Frobenius matrix entries of a MUM operator
  (Beukers--Vlasenko \cite{BVfrob}), which is the conceptual reason a weight-$w$ second
  solution should descend with a factor $p^w$ --- but that is a statement about a $p$-adic
  \emph{limit}, proved for simplicial and hyperoctahedral Calabi--Yau families, and the
  classical Ap\'ery operator is not among them.
\item The one genuine precedent is Gessel's companion congruence, in the form given by
  Beukers \cite[Thm.~1.4]{Beukers2211}: with
  $A'_k=\sum_{m=0}^k\bin km^2\bin{m+k}m^2\bigl(\tfrac1{m+1}+\dots+\tfrac1k\bigr)$ one has
  $A_{kp+\ell}\equiv A_\ell A_k+p\,A'_\ell\,k\,A_k \pmod{p^2}$. This is a
  \emph{weight-one} harmonic companion, and the depth stops at $p^2$. (We have confirmed
  the modulus $p^2$ and the definition of $A'$ against the source; the placement of the
  indices in the $p$-term should be checked against the published version before it is
  relied upon.)
\end{itemize}

So the second row is a genuine gap: weight one is known modulo $p^2$; weight three, the
weight of $\zeta(3)$ and of Ap\'ery's own companion, is not known at all.

\subsection{Results}
\label{ssec:results}

\subsubsection*{(A) The weight-three theorem}
Our first result is elementary and self-contained.

\begin{theorem}[weight-three descent; Theorem~\ref{thm:brow} below]
\label{thm:intro-brow}
Let $p\ge 5$ be prime, let $1\le a<p$ and $0\le r<p$, and put $n=ap+r$. Then $b_a$ and
$p^3b_n$ lie in $\Zp$ and
\[
   p^3\,b_{ap+r}\;\equiv\;b_a\,a_r \pmod p .
\]
\end{theorem}

The proof occupies Section~\ref{sec:weight3}; it needs Lucas' theorem, Kummer's carry
criterion, and nothing else. The three steps that a referee found compressed in the first
write-up --- the borrow count in $\bin nm$, the product-completion of the summation
region, and the interchange defining the tail sums --- are here isolated as
Lemma~\ref{lem:borrow}, Lemma~\ref{lem:fibre} and Lemma~\ref{lem:interchange}
respectively, and proved in full.

\subsubsection*{(B) The master depth law}
Numerically the congruence of Theorem~\ref{thm:intro-brow} is the mod-$p$ shadow of a
sharper and more symmetric statement which we can verify but not prove. Writing
$q=\lfloor n/p\rfloor$, the quantity $\vp(p^3b_na_q-b_qa_n)$ is $\ge 3$, with minimum
exactly $3$, for every prime $5\le p\le 31$ and every $n\le 320$ --- including the
multi-digit range and including $p=5$. Section~\ref{sec:depth} states this as
Conjecture~\ref{conj:master3} and records the exact sweep parameters, the fine structure
of the sporadic boosts above the floor, and two false statements that the data kills. The
descent depth equals the motivic weight of the Ap\'ery limit; this is what makes the
programme's name for it, \emph{the master depth law}, appropriate.

\subsubsection*{(C) The general $\chi$-twisted theorem}
The classification of Ap\'ery-like recurrences produces fifteen sporadic pairs
$(A(n),B(n))$. Running the same experiment on all fifteen produces the paper's main
discovery: the naive law is \emph{false} for seven of them, and the failure is governed by
a quadratic Dirichlet character.

\begin{conjecture}[master form; Conjecture~\ref{conj:master} below]
\label{conj:intro-master}
For each sporadic pair, with $w$ the weight of its Ap\'ery limit, $\chi$ the quadratic
character attached to it, $p\ge 5$ and $q=\lfloor n/p\rfloor$,
\[
  \vp\bigl(p^{w}B_nA_q-\chi(p)B_qA_n\bigr)\;\ge\;w .
\]
\end{conjecture}

\noindent
\VerifiedR{zero failures} for all fifteen pairs, all primes $5\le p\le 103$, all $n\le400$,
and $n\le 1200$ for $p=5,7,11$ (five base-$p$ digits), in exact rational arithmetic; the
floor is exactly $w$ in every single (sequence, prime) cell.

The mod-$p$, single-digit form of Conjecture~\ref{conj:intro-master} we can prove, in a
generality that covers four of the fifteen families outright:

\begin{theorem}[Theorem~\ref{thm:LB} below]
\label{thm:intro-LB}
Let $p\ge5$ and let $B(n)=\sum_kS(n,k)w(n,k)$, where $S$ is a product of binomials in
linear forms and $w$ is a $\Q$-combination of weight-$w$ monomials in harmonic letters
$\Kchi{r}{x}=\sum_{m\le x}\chi(m)m^{-r}$. Under the five hypotheses
\textup{(H1)--(H5)} of Section~\ref{sec:general},
\[
  p^{w}\,B(ap+r)\;\equiv\;\chi(p)^{e}\,B(a)\,A(r)\pmod p
  \qquad(1\le a<p,\;0\le r<p),
\]
where $e\in\{0,1\}$ is the common number of $\chi$-letters per monomial.
\end{theorem}

The mechanism is a two-line lemma (Lemma~\ref{lem:K}): for a completely multiplicative
$\chi$,
\[
   p^r\Kchi ry \;=\; \chi(p)\,\Kchi r{\lfloor y/p\rfloor}\;+\;p^r\cdot(\text{$p$-integral}),
\]
because the $p$-divisible indices $m=jp$ contribute $\chi(jp)(jp)^{-r}
=\chi(p)p^{-r}\chi(j)j^{-r}$. \emph{The factor $\chi(p)$ in the congruence is exactly the
$\chi(p)$ coming out of the $p$-divisible layer of the harmonic weight.} Correspondingly,
$\chi$ is the character of the sequence's Ap\'ery limit: the limit is a rational multiple
of $\zeta(w)$ when $\chi$ is trivial, and of $L(\chi_D,w)$ when $\chi=\chi_D$. The
correlation is $12$ out of $12$ on the sequences whose limit exists.

\subsubsection*{(D) A minimal closed form}
The general theorem needs a harmonic-monomial weight, and finding one for Ap\'ery's $b_n$
produced a formula we have not been able to locate anywhere:
\begin{equation}\label{eq:minimal-intro}
   b_n\;=\;\sum_{k=0}^n\bin nk^2\bin{n+k}k^2\bigl(2\HH3n-\HH3k\bigr).
\end{equation}
Section~\ref{sec:minimal} proves \eqref{eq:minimal-intro} by six explicit
Zeilberger/Gosper certificates, and proves two further closed forms of the same kind, for
Franel's numbers and for $\sum_k\bin nk^4$. Substituted into Theorem~\ref{thm:intro-LB},
\eqref{eq:minimal-intro} reduces Theorem~\ref{thm:intro-brow} to three lines. It also
removes the hypothesis $p\ge5$ for that instance, since the coefficients $2,-1$ are
integers.

\subsubsection*{(E) Sharpness, positioning, formalisation}
Section~\ref{sec:sharpness} documents the failure of the master law at $p=2,3$, and
locates the obstruction where the proof says it should be: in the denominators of the
weight's coefficients. Section~\ref{sec:discussion} positions
Conjecture~\ref{conj:intro-master} against Beukers--Vlasenko's Conjecture~7.5: the two are
statements about the two diagonal entries of one $2\times2$ Frobenius matrix, neither
implies the other, and the character is invisible in the unit-root formulation.
Section~\ref{sec:lean} reports what is machine-verified: Theorem~\ref{thm:intro-LB},
Lemma~\ref{lem:K}, the two first-row Lucas congruences, and two instances of the
second-row congruence are formalised in Lean~4, sorry-free, on the three standard Mathlib
axioms.

\subsection{Evidence classes}
\label{ssec:evidence}

This paper reports a mixture of proofs and of exact finite computations, and it labels
them. We use exactly three markers, and we never upgrade one to another:

\begin{description}[leftmargin=2.2em,style=nextline]
\item[\Proved] a complete proof is written out here.
\item[\Verified] an exact finite computation, always with its range stated. Exact rational
  or modular arithmetic, no floating point. This is evidence and never proof, however
  large the range.
\item[\Open] neither.
\end{description}

Statements carrying only \Verified\ are called \emph{conjectures}, and are stated in
\texttt{conjecture} environments, with the sweep parameters attached. Where a computation
was run twice, by independent implementations, we say so.

\subsection{Notation}
\label{ssec:notation}

Throughout, $p$ is a prime, $\vp$ is the $p$-adic valuation on $\Q$ normalised by
$\vp(p)=1$, and $\Zp=\{x\in\Q:\vp(x)\ge0\}$ is the local ring; a congruence
$x\equiv y\pmod{p^s}$ between rationals means $\vp(x-y)\ge s$. We write $d_n=\lcm(1,\dots,n)$,
$\HH rx=\sum_{m=1}^x m^{-r}$ (so $\HH r0=0$), and for a completely multiplicative
arithmetic function $\chi$,
\[
   \Kchi ry\;=\;\sum_{m=1}^{y}\frac{\chi(m)}{m^{r}} ,
\]
which for $\chi\equiv1$ is $\HH ry$. For a sporadic pair we write $A(n)$ for the integral
solution and $B(n)$ for the second solution normalised by $B(0)=0$, $B(1)=1$; for the
classical $\zeta(3)$ pair we write $a_n$ and $b_n$, with the relation $b_n=6B(n)$. Digits:
for $n=ap+r$ with $0\le r<p$ we call $a=\lfloor n/p\rfloor$ the high digit and $r$ the low
digit, and similarly $k=bp+s$. Finally $\bin nk=0$ whenever $k>n$ or $k<0$, which we use
without comment.
